% Standalone section -- paste into the weekly report.
% Presents the scope decision on its own: the topology/feature-extraction work in this thesis will
% cover three named trunk vessels of the left coronary system (LM, LAD, LCX) rather than the
% coronary tree in full. Segmentation, dominance annotation and centerline extraction still need the
% whole tree; only the population-level feature analysis (objective 7) is being narrowed.
% Citation keys resolve against thesis/refs.bib. cui2017bifurcation, juan2017bifurcation,
% mansouri2024bifurcation, tommasino2024clap and lawton2022revasc were first cited here, so refs_clean.bib has to be
% regenerated after editing: python scripts/make_refs_clean.py
% Do NOT quote tommasino2024clap's HR 4.47: its CI (3.80-6.70) is not centred on the estimate and
% the full text has not been checked. Direction only.

\subsection{Scope: focusing for a start on the left main, LAD and LCX}

As a start I'm narrowing the vessels whose topology I measure and analyze to three: the \textbf{left main
artery (LM)}, the \textbf{left anterior descending artery (LAD)} and the \textbf{left circumflex
artery (LCX)}, rather than the coronary tree in full.

The three belong together as one structure: the LM and the bifurcation at which it divides into the
LAD and the LCX. That structure is where both the clinical stakes and the geometric evidence are
concentrated.

The clinical stakes follow from what the vessels supply. The LM carries all the blood bound for the
left tree, so a lesion there threatens the territories of both its branches at once. Guidelines
reflect that directly: a left main stenosis counts as significant from 50\%, where every other
coronary segment needs 70\% \cite{lawton2022revasc}. The LAD
supplies the front wall of both ventricles and the septum between them, the LCX the side wall of the
left ventricle \cite{devos2024coronary,obrien2007anatomy}. The LAD is also where rupture-prone plaque
concentrates: optical coherence tomography of all three coronary arteries in 131 patients found
thin-cap fibroatheroma in 18.9\% of 534 plaques, clustered in the proximal LAD
\cite{araki2020plaques}.

The geometric evidence sits at the bifurcation, and that is why the LCX is included: a bifurcation
cannot be measured without both of its branches. The angle between the LAD and the LCX is the
geometric feature with the most consistent clinical evidence. A wider angle has gone with more
severe disease in cohorts of 106 to 313 patients
\cite{cui2017bifurcation,juan2017bifurcation,mansouri2024bifurcation}, and with more cardiac events
over follow-up in 499 \cite{tommasino2024clap}. Simulation studies disagree on whether the angle
alone changes the shear on the wall \cite{shen2026geometry}, which is one reason to measure it
together with the diameters of both branches. The flow evidence is concentrated here as well: the
study that found curvature, rather than the tortuosity index, to track low wall shear stress in 127
patients without coronary artery disease measured the LM and its two branches
\cite{kashyap2022tortuosity}. \Cref{fig:week2-left-variants} shows the three vessels in the
ImageCAS-X data, together with the two connectivity variants that change what ``the LM'' is.

\begin{figure}[htbp]
  \centering
  \includegraphics[width=\textwidth]{content/background/figures/06_left_variants.png}
  \caption{The left coronary system in three ImageCAS-X cases (top: the whole left tree; bottom:
    the first 15\,mm from the ostium). Typically the LM divides into the LAD and LCX (left). In 11
    of 800 cases there is no LM and the LAD and LCX leave the aorta separately (middle); in 196 of
    800 a ramus intermedius (IM) makes the LM end in three branches rather than two (right).}
  \label{fig:week2-left-variants}
\end{figure}

These three vessels are where the evidence is, not where all the disease is. The same plaque mapping
found further peaks in the right coronary artery (RCA) \cite{araki2020plaques}, and the field's own
review counts the concentration of studies on the main bifurcation regions among its limitations
\cite{shen2026geometry}. The rest of the tree is not disease-free, and I may return to including
more of it later.

